\section{Why the Frozen First Layer is Crucial: Preventing Divergence}
\label{sec:frozen-layer-crucial}

This section explains why the frozen first layer is not just a simplification, but the \textbf{key mechanism} that prevents divergence and enables the generalization to arbitrary depth $L$.

\subsection{The Problem: What If the First Layer Could Diverge?}

If the first layer were trainable and could diverge, we would face a fundamental problem:

\begin{enumerate}
\item \textbf{Error propagation}: If $w_1(t)$ diverges (e.g., $|w_1(t)| \to \infty$), then:
   \begin{align*}
   H_2(t, c_2; X, W) &= \sum_{k=1}^r L_{c_2,k}\,\mathbb{E}_{C_1}[w_1(t, C_1, k)\,\varphi_1(\langle f_1(C_1), X\rangle)]\\
   &\to \infty \quad \text{(if $w_1$ diverges)}
   \end{align*}
   This would cause all subsequent layers to diverge.

\item \textbf{Support collapse}: If $w_1(t)$ moves toward a degenerate distribution (e.g., all neurons collapse to the same direction), then:
   \[
   \text{supp}(w_1(t, C_1)) \subsetneq \mathbb{R}^d,
   \]
   and universal approximation would fail. This is exactly what happens in full-width networks with i.i.d. initialization for $L \geq 3$ (Corollary 29 of \cite{nguyen2023rigorousframeworkmeanfield}).

\item \textbf{Need for correlated initialization}: To prevent divergence and support collapse in a trainable first layer, we would need:
   \begin{itemize}
   \item Correlated initialization: $(w_1^0(C_1), w_2^0(C_1, \cdot))$ must have full support
   \item Function-dependent design: The correlation structure must be chosen based on the target function
   \item This is exactly what limits full-width proofs to $L=3$ layers
   \end{itemize}
\end{enumerate}

\subsection{The Solution: Frozen First Layer}

\textbf{The frozen first layer solves all three problems simultaneously:}

\begin{enumerate}
\item \textbf{No divergence possible}: Since $f_1(C_1)$ is fixed and independent of $W(t)$:
   \[
   H_1(t, c_1; X, W) = \langle f_1(c_1), X\rangle \quad \text{(constant in $t$)}
   \]
   This cannot diverge because it doesn't depend on $W(t)$ at all.

\item \textbf{Full support maintained automatically}: Since $f_1(C_1)$ are frozen random features (e.g., Gaussian):
   \[
   \text{supp}(f_1(C_1)) = \mathbb{R}^d \quad \text{for all $t \geq 0$}
   \]
   This is preserved \textbf{trivially} because $f_1$ never changes.

\item \textbf{No correlated initialization needed}: Since $f_1(C_1)$ has full support by construction (e.g., Gaussian random features), we can initialize $w_1^0(C_1, k)$ and $w_2^0(C_2)$ independently with standard i.i.d. distributions.
\end{enumerate}

\subsection{Why This Enables Generalization to Arbitrary $L$}

\textbf{The key insight is that we're not actually "stacking" layers in a way that compounds errors.} Instead:

\begin{enumerate}
\item \textbf{Stable base}: The frozen first layer provides a stable, non-diverging base:
   \[
   H_1(t, c_1; X, W) = \langle f_1(c_1), X\rangle \quad \text{(independent of $W(t)$)}
   \]
   This is the \textbf{base case} of our recursion.

\item \textbf{Same procedure at each layer}: For layers $i = 2, \ldots, L$, we apply the \textbf{same bounded procedure}:
   \begin{itemize}
   \item Forward: $H_i$ depends on $H_{i-1}$ through bounded mixing matrix $L^{(i)}$ with $\|L^{(i)}\|_{\infty,1} \le rK$
   \item Backward: $B_k^{(i)}$ depends on $B_k^{(i+1)}$ through bounded operations
   \item All errors bounded in terms of global distance $\mathscr{D}_t(W, \tilde{W})$, not layer-wise errors
   \end{itemize}

\item \textbf{No error compounding}: Because we bound everything in terms of the global distance $\mathscr{D}_t$, errors don't compound exponentially. Each layer contributes a factor $(1+rK)$ to the Gronwall constant, and these multiply (not compound):
   \[
   \text{Gronwall constant} = K_T(1+rK)^{L-2} \quad \text{(polynomial, not exponential)}
   \]
\end{enumerate}

\subsection{The Recursive Structure}

\textbf{We're reusing the same procedure, but with a crucial difference:}

For $L=3$ layers:
\begin{align*}
H_1 &= \langle f_1(C_1), X\rangle \quad \text{(frozen, stable)}\\
H_2 &= \sum_{k=1}^r L_{c_2,k}\,\mathbb{E}_{C_1}[w_1(t, C_1, k)\,\varphi_1(H_1)]\\
\hat{y} &= \mathbb{E}_{C_2}[w_2(t, C_2)\,\varphi_2(H_2)]
\end{align*}

For $L$ layers:
\begin{align*}
H_1 &= \langle f_1(C_1), X\rangle \quad \text{(frozen, stable - SAME)}\\
H_2 &= \sum_{k=1}^r L^{(2)}_{c_2,k}\,\mathbb{E}_{C_1}[w_1(t, C_1, k)\,\varphi_1(H_1)] \quad \text{(SAME procedure)}\\
H_3 &= \sum_{k=1}^r L^{(3)}_{c_3,k}\,\mathbb{E}_{C_2}[w_2(t, C_2, k)\,\varphi_2(H_2)] \quad \text{(SAME procedure)}\\
&\vdots\\
H_L &= \mathbb{E}_{C_{L-1}}[w_{L-1}(t, C_{L-1}, k)\,\varphi_{L-1}(H_{L-1})] \quad \text{(SAME procedure)}
\end{align*}

\textbf{The key is that:}
\begin{itemize}
\item The base case $H_1$ is frozen and stable (can't diverge)
\item Each subsequent layer applies the same bounded procedure
\item All errors are bounded in terms of the global distance $\mathscr{D}_t$
\item The Gronwall constant accumulates multiplicatively: $(1+rK)^{L-2}$
\end{itemize}

\subsection{What Would Happen With a Trainable First Layer?}

If we tried to make the first layer trainable, we would need to:

\begin{enumerate}
\item \textbf{Prove $w_1(t)$ doesn't diverge}: This requires showing that the gradient flow keeps $w_1(t)$ bounded. This is non-trivial and requires additional assumptions.

\item \textbf{Prove support remains full}: We need $\text{supp}(w_1(t, C_1)) = \mathbb{R}^d$ for all $t$. This is exactly what requires correlated initialization in the full-width case.

\item \textbf{Handle error compounding}: If $w_1(t)$ could move, then errors in $w_1$ would propagate to $H_2$, then to $H_3$, etc., potentially compounding exponentially.

\item \textbf{Require function-dependent initialization}: The correlation structure $(w_1^0(C_1), w_2^0(C_1, \cdot))$ would need to be chosen based on the target function, limiting applicability.
\end{enumerate}

\textbf{This is exactly why full-width proofs are limited to $L=3$ layers!}

\subsection{The Key Insight: Frozen = Stable Base}

\textbf{The frozen first layer is not a limitation—it's the key that unlocks arbitrary depth:}

\begin{theorem}[Frozen first layer enables arbitrary depth]
If the first layer is frozen (i.e., $f_1(C_1)$ is independent of $W(t)$), then:
\begin{enumerate}
\item The base case $H_1 = \langle f_1(C_1), X\rangle$ is stable and cannot diverge
\item Full support $\text{supp}(f_1(C_1)) = \mathbb{R}^d$ is maintained trivially
\item Standard i.i.d. initialization suffices (no correlated initialization needed)
\item The same bounded procedure can be applied at each subsequent layer
\item The Gronwall constant grows polynomially: $(1+rK)^{L-2}$
\end{enumerate}
Therefore, global convergence holds for arbitrary depth $L$ with standard initialization.
\end{theorem}

\subsection{Conclusion}

\textbf{To answer the question: "How do we avoid divergence when stacking layers?"}

The answer is: \textbf{We don't stack layers in a way that compounds errors. Instead:}

\begin{enumerate}
\item \textbf{Frozen first layer provides stable base}: $H_1 = \langle f_1(C_1), X\rangle$ is independent of $W(t)$, so it cannot diverge.

\item \textbf{Same bounded procedure at each layer}: Each layer $i = 2, \ldots, L$ applies the same procedure with bounded mixing matrix $L^{(i)}$ satisfying $\|L^{(i)}\|_{\infty,1} \le rK$.

\item \textbf{Global distance bounding}: All errors are bounded in terms of the global distance $\mathscr{D}_t(W, \tilde{W})$, not layer-wise errors that would compound.

\item \textbf{Polynomial growth}: The Gronwall constant grows as $(1+rK)^{L-2}$, which is polynomial in $L$ (not exponential).

\item \textbf{Reusing the same procedure}: We're not inventing new techniques for each layer—we're applying the same bounded procedure recursively, starting from a stable frozen base.
\end{enumerate}

\textbf{The frozen first layer is the key that makes everything work.} Without it, we'd be back to needing correlated initialization and would be limited to $L=3$ layers, just like the full-width case.
